\section{Error Measures and Basic Bounds}

Fix the format $(E,M)$ of Section 3. In this section write
$\mathbb F=\mathbb F_{E,M}$ and $\rnd=\rnd_{E,M}$, and let
$F_{\max}$, $e_{\min}$, and $e_{\max}$ retain their earlier meanings.
For $x\in\R$, call $x$ in the \emph{normal range} when
$2^{e_{\min}}\le |x|\le F_{\max}$ and in the \emph{subnormal range}
when $|x|<2^{e_{\min}}$.

\begin{definition}[Absolute and Relative Errors]
For an exact value $x\in\R$ and an approximation $\widehat{x}\in\mathbb F$,
set
\[
\begin{aligned}
 \operatorname{err}_{\mathrm{abs}}(x,\widehat{x})
   &\coloneqq |x-\widehat{x}|,\\
 \operatorname{err}_{\mathrm{rel}}(x,\widehat{x})
   &\coloneqq \frac{|x-\widehat{x}|}{|x|}\qquad(x\ne0).
\end{aligned}
\]
\end{definition}

\begin{definition}[Unit in the Last Place]
The \emph{unit in the last place} is the function
$\ulp\colon\R\to\R_{>0}$ defined by
\[
\ulp(x)\coloneqq
\begin{cases}
 2^{e-M},&2^e\le|x|<2^{e+1}
       \text{ for some integer }e\ge e_{\min},\\
 2^{e_{\min}-M},&|x|<2^{e_{\min}}.
\end{cases}
\]
The \emph{ULP error} of $\widehat{x}\in\mathbb F$ is
\[
 \operatorname{err}_{\mathrm{ulp}}(x,\widehat{x})
   \coloneqq\frac{|x-\widehat{x}|}{\ulp(x)}.
\]
Thus ULP error measures the discrepancy in units of the local
floating-point spacing (within a common binade). The
\emph{machine epsilon} (unit roundoff) of the format is
\[
 \varepsilon\coloneqq\tfrac12\ulp(1)=2^{-M-1}.
\]
\end{definition}

\begin{proposition}[Relative and ULP Errors]
Let $x\in\R\setminus\{0\}$ and $\widehat{x}\in\mathbb F$.
Then
\[
 \operatorname{err}_{\mathrm{ulp}}(x,\widehat{x})
 \le\frac{\operatorname{err}_{\mathrm{rel}}(x,\widehat{x})}
 {\varepsilon}.
\]
If $x$ is in the normal range, then also
\[
 \operatorname{err}_{\mathrm{ulp}}(x,\widehat{x})
 \ge\frac{\operatorname{err}_{\mathrm{rel}}(x,\widehat{x})}
 {2\varepsilon}.
\]
\end{proposition}
\begin{proof}
The definitions give, without dividing by an error that might vanish,
\[
 \operatorname{err}_{\mathrm{ulp}}(x,\widehat{x})
 =\operatorname{err}_{\mathrm{rel}}(x,\widehat{x})
   \frac{|x|}{\ulp(x)}.
\]
For $2^e\le|x|<2^{e+1}$ with $e\ge e_{\min}$,
\[
 2^M\le\frac{|x|}{\ulp(x)}<2^{M+1}=\frac1\varepsilon.
\]
In the subnormal range, the same ratio is at most $2^M$;
this proves the first inequality. In the normal range its lower
bound is $2^M=1/(2\varepsilon)$, proving the second.

The normal-range hypothesis cannot be dropped. Indeed, for
$x=2^{e_{\min}-M}$ and $\widehat{x}=0$, both the relative and ULP
errors equal $1$, but $1/(2\varepsilon)=2^M>1$.
\end{proof}

\begin{proposition}[Rounding Error Bounds]
For round-to-nearest, ties-to-even, the following estimates hold:
\begin{enumerate}[label=(\roman*)]
 \item If $x$ is in the normal range, then
 \[
 \operatorname{err}_{\mathrm{rel}}(x,\rnd(x))\le\varepsilon.
 \]
 \item If $x$ is in the subnormal range, then
 \[
 \operatorname{err}_{\mathrm{abs}}(x,\rnd(x))
 \le 2^{e_{\min}}\varepsilon.
 \]
 \item If $|x|\le F_{\max}$, then
 \[
 \operatorname{err}_{\mathrm{ulp}}(x,\rnd(x))\le\tfrac12.
 \]
\end{enumerate}
\end{proposition}
\begin{proof}
For $2^e\le|x|<2^{e+1}$ in the normal range, consecutive
floating-point values are separated by at most $2^{e-M}$ locally.
Thus nearest rounding gives
\[
 \frac{|x-\rnd(x)|}{|x|}
 \le\frac{2^{e-M-1}}{2^e}=\varepsilon.
\]
In the subnormal range the spacing is $2^{e_{\min}-M}$, so
\[
 |x-\rnd(x)|\le\tfrac12 2^{e_{\min}-M}
 =2^{e_{\min}}\varepsilon.
\]
Finally, for $|x|\le F_{\max}$, nearest rounding is within half
the local spacing $\ulp(x)$, giving (iii).
\end{proof}

\begin{definition}[Underflow and Overflow]
For an exact nonzero arithmetic result $z$, the \emph{underflow range}
is $0<|z|<2^{e_{\min}}$. A result beyond the finite range satisfies
$|z|>F_{\max}$; under round-to-nearest, actual \emph{overflow to infinity}
occurs when $|z|\ge K_{E,M}$. These two upper thresholds need not
coincide.
\end{definition}

\begin{theorem}[Error of Floating-Point Operations]
Let $a,b\in\mathbb F$ and choose corresponding operations
\[
 \ast\in\{+,-,\times,\div\},\qquad
 \circledast\in\{\oplus,\ominus,\otimes,\odiv\},
\]
with $b\ne0$ for division. If $z=a\ast b$ satisfies
$|z|\le F_{\max}$, then
\[
 a\circledast b=z(1+\delta)+\delta',\qquad
 |\delta|\le\varepsilon,\quad
 |\delta'|\le2^{e_{\min}}\varepsilon,
\]
where one may choose $\delta'=0$ in the normal range and
$\delta=0$ in the subnormal range (including $z=0$).
For addition and subtraction, one may always choose $\delta'=0$.
\end{theorem}
\begin{proof}
By definition, $a\circledast b=\rnd(z)$. In the normal range,
the relative rounding bound yields
$\rnd(z)=z(1+\delta)$ with $|\delta|\le\varepsilon$.
In the subnormal range, the absolute bound yields
$\rnd(z)=z+\delta'$ with
$|\delta'|\le2^{e_{\min}}\varepsilon$.

For $\ast\in\{+,-\}$, every element of $\mathbb F$ is an integer
multiple of $2^{e_{\min}-M}$, so $a\pm b$ is also such a multiple.
If this exact sum or difference is in the subnormal range, it is
representable and rounding is exact. Thus $\delta'=0$ in both ranges.
\end{proof}

\begin{example}[Error in Repeated Multiplication]
Let $a,b,c\in\mathbb F\setminus\{0\}$, and suppose both intermediate
exact products are in the normal range and do not overflow.
By the preceding theorem, for some $|\delta_1|,|\delta_2|\le\varepsilon$,
\[
\begin{aligned}
 (a\otimes b)\otimes c
 &=\bigl(ab(1+\delta_1)c\bigr)(1+\delta_2)\\
 &=abc(1+\delta_1)(1+\delta_2).
\end{aligned}
\]
Consequently, with $x=abc$ and $\widehat{x}=(a\otimes b)\otimes c$,
\[
\begin{aligned}
 \operatorname{err}_{\mathrm{rel}}(x,\widehat{x})
 &\le |\delta_1|+|\delta_2|+|\delta_1\delta_2|
 \le 2\varepsilon+\varepsilon^2,\\
 \operatorname{err}_{\mathrm{ulp}}(x,\widehat{x})
 &\le 2+\varepsilon.
\end{aligned}
\]
To first order the relative error is at most $2\varepsilon$.
\end{example}

\begin{example}[Error in Repeated Addition]
Let $a,b,c\in\mathbb F$ and set $s=a+b$, $u=a+b+c\ne0$.
Assume $|s|, |(a\oplus b)+c|\le F_{\max}$. For rounded addition,
the theorem gives $\delta'=0$ even for subnormal results; hence
\[
\begin{aligned}
 (a\oplus b)\oplus c
 &=\bigl((a+b)(1+\delta_1)+c\bigr)(1+\delta_2)\\
 &=u(1+\delta_2)+s\delta_1(1+\delta_2),
 \qquad |\delta_1|,|\delta_2|\le\varepsilon.
\end{aligned}
\]
Writing $x=u$ and $\widehat{x}=(a\oplus b)\oplus c$, we obtain
\[
\begin{aligned}
 \operatorname{err}_{\mathrm{rel}}(x,\widehat{x})
 &\le\varepsilon\left(1+\frac{|a+b|}{|a+b+c|}\right)
   +\varepsilon^2\frac{|a+b|}{|a+b+c|},\\
 \operatorname{err}_{\mathrm{ulp}}(x,\widehat{x})
 &\le 1+\frac{|a+b|}{|a+b+c|}
       +\varepsilon\frac{|a+b|}{|a+b+c|}.
\end{aligned}
\]
The ratio $|a+b|/|a+b+c|$ can amplify the first rounding error.

For example, when $2^{M+1}\in\mathbb F$, take
$a=1$, $b=2^{M+1}$, and $c=-2^{M+1}$. Ties-to-even yields
\[
 (a\oplus b)\oplus c=0,
 \qquad a+b+c=1,
\]
with relative error $1$ and ULP error $2^M$.
\end{example}
